% Generated draft: requires mathematical and source review.
\documentclass[11pt]{article}
\usepackage[margin=25mm]{geometry}
\usepackage{fontspec}
\setmainfont{Latin Modern Roman}
\usepackage{amsmath,amssymb,mathtools,unicode-math}
\setmathfont{Latin Modern Math}
\usepackage{enumitem,needspace,xurl,hyperref}
\setlength{\parindent}{0pt}
\setlength{\parskip}{5pt}
\setlength{\emergencystretch}{3em}
\newcommand{\N}{\mathbb{N}}
\newcommand{\Z}{\mathbb{Z}}
\newcommand{\Q}{\mathbb{Q}}
\newcommand{\R}{\mathbb{R}}
\newcommand{\C}{\mathbb{C}}
\newcommand{\F}{\mathbb{F}}
\newcommand{\A}{\mathbb{A}}
\newcommand{\PP}{\mathbb{P}}
\newcommand{\E}{\mathbb{E}}
\newcommand{\eps}{\varepsilon}
\newcommand{\dd}{\,\mathrm{d}}
\begin{document}
\section*{UnsolvedMath ID: 701355}
\textbf{Source problem code:} catalogue entry
\subsection{Definitions and mathematical statement}
The source concerns finite groups. It does not specify the coefficient field or the representation convention. The formulation below uses finite-dimensional complex representations, subject to confirmation that this is the intended convention.

For a finite group \(G\), let \(\operatorname{Irr}_{\C}(G)\) denote the set of isomorphism classes of irreducible finite-dimensional complex representations of \(G\). Under this convention, call \(G\) \emph{simply reducible}, or an \emph{SR-group}, if both of the following hold:
\begin{enumerate}
\item For every \(g\in G\), there exists \(x\in G\) such that \(xgx^{-1}=g^{-1}\).
\item For every \(V,W\in\operatorname{Irr}_{\C}(G)\), the tensor product with diagonal \(G\)-action satisfies
\[
V\otimes_{\C}W\cong\bigoplus_{U\in\operatorname{Irr}_{\C}(G)}U^{\oplus m_U(V,W)},
\qquad m_U(V,W)\in\{0,1\}.
\]
\end{enumerate}
Here the two representations \(V\) and \(W\) may be equal.

Define \(G^{(0)}=G\) and \(G^{(i+1)}=[G^{(i)},G^{(i)}]\), where \([H,H]\) is the subgroup generated by the commutators of elements of \(H\). A group is \emph{soluble} if \(G^{(d)}=1\) for some integer \(d\geq0\); its \emph{derived length} \(\operatorname{dl}(G)\) is the least such \(d\).

For a finite soluble group, interpret its \emph{nilpotent length} as its Fitting height \(h(G)\): the least integer \(r\geq0\) for which there exists a series
\[
1=N_0\leq N_1\leq\cdots\leq N_r=G,
\qquad N_i\trianglelefteq G,
\]
with every factor \(N_i/N_{i-1}\) nilpotent. A group \(H\) is nilpotent if its lower central series, defined by \(\gamma_1(H)=H\) and \(\gamma_{j+1}(H)=[\gamma_j(H),H]\), reaches \(1\). Thus \(h(1)=\operatorname{dl}(1)=0\).

L. S. Kazarin's questions are:
\begin{enumerate}
\item[(a)] Does every finite soluble SR-group \(G\) satisfy \(h(G)\leq5\)?
\item[(b)] Does there exist an integer \(c\geq0\), independent of \(G\), such that every finite soluble SR-group \(G\) satisfies \(\operatorname{dl}(G)\leq c\)?
\end{enumerate}
These questions use the complex-representation convention stated above; the omitted coefficient field in the source remains an unresolved qualification.
\subsection{Short English statement}
Kazarin asks whether every finite soluble simply reducible group has nilpotent length at most \(5\), and whether the derived lengths of all such groups admit one uniform constant bound. Simply reducible means that every element is conjugate to its inverse and every tensor product of two irreducible representations is multiplicity-free. The source omits the representation field; this draft explicitly uses \(\C\), pending confirmation.
\subsection{Sources}
\begin{itemize}
\item Local UnsolvedMath registry: supplied complete source record, permanent ID \(701355\), exact source problem code \texttt{catalogue entry}, title \emph{Kourovka Notebook Problem 17.53}. Supplied upstream revision: \texttt{d371ee917ac9647fbb06cd6fc529deec6f3dda82}; record updated \texttt{2026-10-16T14:00:00Z}.
\item Supplied source attribution: \emph{The Kourovka Notebook (Unsolved Problems in Group Theory)}, original problem 17.53, 17th issue (2010); also identified in the supplied background with the 21st issue edition (2026), \texttt{arXiv:1401.0300v48}. The supplied collection description names E. I. Khukhro and V. D. Mazurov as editors. These bibliographic details have not been independently verified.
\item Proposed by L. S. Kazarin in 2010, according to the supplied record. Its status is labelled \emph{open}; this is supplied registry metadata, not an independently verified assessment of the literature. No separate literature-verification date or literature-source links were supplied.
\end{itemize}
\end{document}
